\documentclass[twoside]{article}
\usepackage[a4paper,margin=0.5in]{geometry}
\usepackage{gregosheet}
\usepackage{gregosheet-psalm}
\input{styles}
\input{ordo-cantus-officii-sheets}

\title{Nagyhét 2026}
\subtitle{A szentendrei ferences kápolna szertartásrendje}

\begin{document}

\maketitle

\section{Virágvasárnap, az Úr szenvedésének vasárnapja}

\commentary{Ezen a napon ünnepli az Egyház Krisztus Urunk Jeruzsálembe való bevonulását a húsvéti %
misztérium beteljesítésére.}

\subsection{Körmenet}

\subsubsection{I. stáció: Az ágak megáldása}

\commentary{A pap, a diakónus és a segédkezők már a szentmiséhez felöltözve a kápolna végébe %
vonul. Eközben a következő antifónát énekeljük:}

\parttitle[Hosanna Filio David]{Anttifóna}

\begin{gregosheet}[tone=7. lá tónus]
  \melody{<ô-1--1t---5--:-5--3---4--5z-5--:-5zZ5-4-:-4-4--4--4--4--2--3---1-,-0--2--4--4---4--4t-rR3-:-1--1t---5x-4-2--3---1---1--,,-3-4-[-…-[-4-5-6-†-4-,-[-6-5-4-†-rR3-2-.}
  \lyrics{   Ho-zsan-na * Dá-vid Fi-á--nak, ál-dott, a-ki az Úr ne-vé-ben jő, Iz-ra-el-nek Ki-rá-lya,  Ho-zsan-na a ma-gas-ság-ban.}
\end{gregosheet}

\commentary{Ha az idő engedi, a 117. zsoltár verseivel énekeljük:}

\begin{psalmsheet}[tone=7, continuous]
  \psalmtext{
    1. Hálát adjatok az Úrnak, mert jó, * mert az ő irgalmassága örökké való. \\
    2. A Kő, amelyet megvetettek az építők, * az lett most Szegletkővé. ANT. \\
    3. Az Úrnak műve ez * és csodálatos a mi szemeink előtt. \\
    4. Tartsatok ünnepeket lombos ágakkal * föl, egészen a magas oltárig. ANT.
}
\end{psalmsheet}

\commentary{A pap a szokott módon üdvözli a népet (ahogyan a mise elején szokta); rövid intelmet %
ad, hogy a hívek tevékenyen és tudatosan vegyenek részt ennek a napnak a megünneplésében. Ezekkel %
vagy hasonló szavakkal fordul hozzájuk:}

\p Kedves Testvéreim! A Nagyböjt eleje óta szívünket bűnbánattal és jócselekedetekkel már
előkészítettük. Ma azért gyűltünk össze, hogy az egész Egyházzal együtt lélekben előre átéljük
Urunk húsvéti misztériumát vagyis szenvedését és feltámadását. Ő ugyanis ezt a misztériumot akarta
beteljesíteni, amikor fölment városába, Jeruzsálembe. Szívünk élő hitével és odaadásával emlékezünk
meg erről az üdvösségszerző bevonulásról. Lépjünk mi is az Úr nyomába, hogy amint kegyelméből
társai vagyunk a kereszthordozásban, ugyanúgy feltámadásának és örök életének is részesei lehessünk.

\parttitle{Könyörgés}

\commentary{Az intelem után a pap a következő két könyörgés egyikét imádkozza összetett kézzel:}

\p Könyörögjünk! Mindenható, örök Isten, szenteld meg áldásoddal \textcolor{red}{†} ezeket az ágakat,
hogy mi, akik most örvendezve lépünk Krisztus Király nyomába, őáltala egykor az örök Jeruzsálembe
is eljussunk. Aki él és uralkodik mindörökkön-örökké.

\h Ámen.

\commentary{Vagy:}

\p Könyörögjünk! Urunk, Istenünk, növeld azoknak hitét, akik benned remélnek. Hallgasd meg könyörgő
imánkat, hogy mi, akik ma virágzó ágakat hozunk a diadalmasan bevonuló Krisztus Király elé, a
jócselekedetek gyümölcsét teremjük őáltala. Aki él és uralkodik mindörökkön-örökké.

\h Ámen.

\commentary{A pap szenteltvízzel meghinti az ágakat, minden kísérő szöveg nélkül.}

\parttitle[Mt 21,1-11]{Evangélium}
\p Az Úr legyen veletek!

\h És a te l\underline{e}lkeddel!

\p Evangélium Szent Mát\underline{é} könyvéből.

\h Dicsőség nek\underline{e}d, Istenünk!

\p Amikor Jeruzsálemhez közeledve az Olajfák-hegyére, Betfagéba értek, Jézus elküldte két tanítványát
ezekkel a szavakkal: „Menjetek előre a szemközti faluba. Ott mindjárt találni fogtok egy szamarat
megkötve és vele csikóját. Oldjátok el és vezessétek hozzám! Ha valaki szólna valamit, mondjátok,
hogy az Úrnak van rá szüksége, és mindjárt elengedi őket.”

Ez azért történt, hogy beteljesedjék, amit a próféta jövendölt: „Mondjátok meg Sion
lányának: Íme, a királyod érkezik hozzád, szelíden, szamárháton ülve, egy teherhordó
állat csikóján.”

A tanítványok elmentek s úgy tettek, ahogy Jézus meghagyta nekik. Elhozták a szamarat és a
csikóját, letakarták ruháikkal, ő pedig felült rá.

A tömegből nagyon sokan az útra terítették ruháikat, mások ágakat tördeltek a fákról
és az útra szórták. Az előtte járó és az utána vonuló tömeg így kiáltott: „Hozsanna Dávid fiának! Áldott, aki az Úr nevében jön! Hozsanna a magasságban!”

Amikor beért Jeruzsálembe, megmozdult az egész város, és kérdezgették: „Ki ez?” A
tömeg ezt felelte: „Ő a Próféta, Jézus, a galileai Názáretből.”

\p Ezek az evangéli\underline{u}m igéi.

\h Áldunk tég\underline{e}d, Krisztus!

\parttitle{Homília}

\commentary{Az evangélium után, ha alkalmas, a pap rövid homíliát mondhat. A körmenet elindulására
vagy maga a pap, vagy más erre alkalmas segédkező adja meg a jelt, ezekkel vagy hasonló szavakkal:}
\p Kedves Testvéreim! Kövessük a Jézust ünneplő jeruzsálemi nép példáját, és induljunk békével!

\h Krisztus nevében, ámen.

\commentary{A körmenet elindul a szentély felé.}

\subsubsection{II. stáció: Az ágak leterítése}

\commentary{A körmenet egy alkalmas helyen megáll, a pap átveszi a körmeneti keresztet, miközben
gyermekek lépnek elébe, ágakat emelnek a magasba, majd leterítik azokat a körmenet elé a földre.
A pap a kereszttel kezében áthalad az ágak fölött. Eközben az előénekesek a következő antifónát
intonáják.}

\parttitle[Pueri hebraeorum portantes]{Antifóna}

\begin{gregosheet}
  \melody{<X-1--3--1---0---3----4---3t-5--:-5---5--5u--5---5-4--3--4t--5-:-5-4--5--3----4--eE2-1--,-3----eE2-1---0--0---1---3r--3--,-3--4t---tT4-2-3--2---1---1-.}
  \lyrics{   Je-ru-zsá-lem gyer-mek né-pe * pál-ma-fák-nak á-ga-it hoz-ván e-lé-be ment az Úr--nak, s_ki-ál--tot-ta han-gos szó-val: Ho-zsan-na  a ma-gas-ság-ban!}
\end{gregosheet}

\subsubsection{III. stáció: A ruhák leterítése}

\commentary{Gyermekek díszes ruhát hoznak és leterítik az útra. A pap, ismét kereszttel kezében,
áthalad rajtuk. Eközben az előénekesek a következő antifónát intonálják:}

\parttitle[Pueri hebraeorum vestimenta]{Antifóna}

\begin{gregosheet}
  \melody{<X-1--3--1---0---3----4---3t-5--:-5--5---5u-5---5--5--4--3--4t--5-:-5----4---5---3---4--eE2-1---,-3--eE2--qQ0-0--1---3--3r--3-,-3--4-----5--tT4-2--3--2--1--1-.}
  \lyrics{   Je-ru-zsá-lem gyer-mek né-pe * le-ter-ít-vén ru-há-it az út-ra,  nagy szó-val így ki-ál--tott: Ho-zsan-na  Dá-vid fi-á-nak,  ál-dott, ki jő  az Úr ne-vé-ben!}
\end{gregosheet}

\subsubsection{IV. stáció: Hódolat Krisztus Király előtt}

\commentary{A körmenet az szentély lépcsőjénél megáll, a pap a keresztet kezébe veszi, s a kereszt
elé járulnak a verzusok énekesei (lehetőség szerint gyerekek). Az éneket ők kezdik, s mindnyájan
megismételjük a refrént. Ezután éneklik a verseket. A versek után közösen megismételjük a refrént.}

\parttitle[Gloria, laus]{Himnusz}

\begin{gregosheet}
  \melody{<X-5--rR3-4---3--4--4---5---3--2--1--:-3---3---1--3--3--2----,-3t-5---5----4--3--4----5--5x--3---2----1-0---eE2-1--1-.}
  \lyrics{   Di-cső-ség és di-csé-ret Te-né-ked, Meg-vál-tó Ki-rá-lyunk, ki-nek gyer-me-ki szép se-reg mon-dott é-kes é---ne-ket.}
\end{gregosheet}

\commentary{Verzusok:}

\begin{gregosheet}
  \melody{<X-1--[----------5u-4--,-7--7---8---7---™6--5-4----5---5-,-[-----------5u-4-,-7----7---8-----7----™6-5---4--5--5-----,,-1-[-5u-4-,-7-7-8-7-™6-5-4-5-5-,-[-5u-4-,-7-7-8-7-™6-5-4-5-5-,,-1-[-5u-4-,-7-7-8-7-™6-5-4-5-5-,-[-5u-4-,-7-7-8-7-™6-5-4-5-5-,,-1-[-5u-4-,-7-7-8-7-™6-5-4-5-5-,-[-5u-4-,-7-7-8-7-™6-5-4-5-5-,,-1-5-5-5-5-5u-4-,-7-7-8-7-™6-5-4-5-5-,-[-5u-4-,-7-7-8-7-™6-5-4-5-5-.}
  \lyrics{   Iz-raelnek_Ki-rá-lya: Dá-vid-nak rég-től i-gért sar-ja, ki_az_Úr_ne-vé-ben jösz-hoz-zánk, mind-ö--rök-ké Ál-dott! REF. An-gyalseregek_a menny-ben mind-nyá-jan ma-gasz-tal-nak Té-ged, s_velük_a_halandó em-ber, és min-den te-remt-mé-nyed di-csér. REF. A zsidóknak né-pe pál-mák-kal jött az ú-ton e-léd, mi_is_könyörgé-sek-kel el-jöt-tünk í-me é-nek szó-val. REF. Mi-dőn_szenvedni men-tél, hó-dol-va di-csér-tek ők Té-ged, most,_midőn_mennyben tró-nolsz, mi zeng-jük e-lőt-ted a him-nuszt. REF. Tet-szett se-re-gük né-ked, tes-sék ma ne-ked szen-telt né-ped, mert_kegyelmes_ki-rály vagy, s_min-den jót tet-szé-sed-del fo-gadsz. REF.}
\end{gregosheet}

\commentary{A pap az oltárhoz megy, a szokott módon tiszteletadást végez, és esetleg megtömjénezi.
Majd minden más bevezető részt elhagyva, a Könyörgéssel kezdi a misét, és a szokott módon folytatja.}

\subsection{Mise}

\parttitle{Könyörgés}

\p Mindenható, örök Isten, Üdvözítőnk példát adott nekünk az alázatosságra, amikor szent
akaratodból emberi testet öltött és vállalta a kereszthalált. Segíts jóságosan, hogy
kínszenvedésének tanítását megértsük, és társai lehessünk feltámadásában. Aki veled él és uralkodik
a Szentlélekkel egységben, Isten mindörökkön-örökké.

\h Ámen.

\parttitle[Iz 50,4-7]{Olvasmány}

\lec Olvamány Izajás próféta könyvéből

Izajás így jövendölt a Megváltóról: Ezt mondja az Úr Szolgája: „Isten, az Úr tanította nyelvemet,
hogy az igével támasza lehessek a megfáradtaknak. Reggelenként ő teszi figyelmessé fülemet, hogy rá
hallgassak, mint a tanítványok. Isten, az Úr nyitotta meg fülemet. S én nem álltam ellen és nem
hátráltam meg. Hátamat odaadtam azoknak, akik vertek, arcomat meg azoknak, akik tépáztak. Nem
rejtettem el arcomat azok elől, akik gyaláztak és leköpdöstek. Isten, az Úr megsegít, ezért nem
vallok szégyent. Arcomat megkeményítem, mint a kőszikla, s tudom, hogy nem kell szégyenkeznem.

\lec Ez az Isten igéje.

\h Istennek legyen hála!

\parttitle[Tenuisti manum]{Graduále}

\begin{gregosheet}
  \melody{<ô-Ÿ---------------2--1---1r-4--:-£---------------------------2-----3r-5--,-£-----------------------2--1-----2r-4-.}
  \lyrics{   Te_tartottad_az én job-bo-mat, és_a_te_akaratod_szerint_ve-zetsz en-gem, és_a_te_dicsőségedbe_be-fo-gadsz en-gem.}
\end{gregosheet}

V. Mily jó Izráel Ist\underline{e}ne † az egyenes szívű\underline{e}khez, * az én lábaim mégis csaknem megtán/tor\underline{o}dtak. TE TARTOTTAD... \\
V. ! M\underline{e}rt megbotránkoztam a bűnösök m\underline{i}att, * látván a bűnösök bé/kess\underline{é}gét. TE TARTOTTAD...

\parttitle[Fil 2,6-11]{Szentlecke}

\lec Szentlecke Szent Pál apostolnak a filippiekhez írt leveléből

Testvéreim! Krisztus Jézus, mint Isten, az Istennel való egyenlőséget nem tartotta olyan dolognak,
amelyhez feltétlenül ragaszkodnia kell, hanem szolgai alakot öltött, kiüresítette önmagát, és
hasonló lett az emberekhez. Megalázta önmagát, és engedelmes lett a halálig, mégpedig a
kereszthalálig. Ezért Isten felmagasztalta őt, és olyan nevet adott neki, amely fölötte van minden
névnek, hogy Jézus nevére hajoljon meg minden térd a mennyben, a földön, és az alvilágban, és
minden nyelv hirdesse az Atyaisten dicsőségére, hogy Jézus Krisztus az Úr.

\lec Ez az Isten igéje.

\h Istennek legyen hála!

\parttitle{Traktus}

\begin{gregosheet}
  \melody{<X-3--4--¥--------------------------:¥----------------7--rR3-:-¥-----------------7--4--uU6-zZ5-4---3r-4--,-¥----------------------------------------7----rR3--:-4--6--6---6--7---4--uU6Z5T4-3r-4-,-©-------------------6----iI7-uJ4R3-,-3--4--6--6---7--uJ4uU66Z5T4-3r-4-,-©------------6---8-7---u7J4R3-,-3---4---6--7---7---4--7--6---zZ5T4-3r-4--,-¥----------------------------7--rR3-3--,-¥-------------------5---4---3---4---ed13r-tT4-,-¥--------------------------------7--rR3-3--,-¥--------------------5-4---3---4--ed13r-tT4-,-5----6---¦---------------------------------6---7--rR3-,-¥--------------------------------5---4--3r--4--,-5z--¦-------------7----6--7---rR3-,-¥------------------5---4--3r-4--,-©----------------------6---iI7-u7J4R3-,-¥---------------------------------:¥-------------uJ4-uU66Z5-4--3r--4--,-©----------------------------6---8---7--u7J4R3-,-.}
  \lyrics{   Is-te-nem,_Istenem,_tekints_reám, miért_hagytál_el en-gem?  Miért_vagy_messze ki-ál-tá--som sza-vá-tól? Én_Istenem,_napestig_kiáltok,_és_meg_nem hall-gatsz? És éj-jel és nem fi-gyelsz  re-ám? Pedig_te_a_szent_he-lyen la-kol,     Iz-ra-el-nek di-cső---------sé-ge! Tebenned_bíz-tak a-tyá-ink,     bíz-tak és meg-sza-ba-dí-tot-tad   ő--ket. Hozzád_kiáltottak,_és_megsza-ba-dul-tak, Tebenned_bíztak,_és meg nem szé-gye-nül---tek.  Én_pedig_féreg_vagyok_már_és_nem is em--ber, emberek_gyalázata_és a nép meg-ve-tett--je.   Mind-nyá-jan,_akik_látnak_engem,_megcsúfol-nak en-gem;  ajkukkal_megszólanak,_és_fejüket haj-to-gat-ják. „Az Úrban_bízott, ment-se meg őt,   szabadítsa_meg,_ha ked-ve-li őt!” Csak_néznek_és_szemlél-nek en--gem,     elosztották_maguk_között_ruháimat, és_köntösömre sor-sot    ve-tet-tek. Szabadíts_meg_engem_az_orosz-lán szá-já-ból,     és a vadak támadásától engem, meg-a-lá - zottat. 9.Akik félitek az Urat, dicsérjé-tek Őt, * Jákobnak ivadékai, dicsőít-sé - tek Őt! 10.Hirdessétek az Urat a jövendő nemzedéknek, hirdessék az egek az Ő i-gazsá - gát a népnek, mely egykor megszü-le - -tik,*me-lyet majdan az Úr te - remt!}
\end{gregosheet}

\parttitle{Passió}
\end{document}
